\section{Conclusion}

Agentic robotics has potential to build on rapidly-advancing LLM models to provide a bridge between Good Old Fashioned Engineering (GOFE) and model-free VLA policies.  As shown in the Ablation experiments, GaP's graph-structured agentic robotics facilitates modularity, multi-agent integration, and self-learning.  As robot applications become more complex and skill libraries grow, more sophisticated agent harnessing will be required.

%In this paper, we introduced the Variational Automation (VA) task class for Robot Learning, formalizing this persistent industrial challenge and providing 8 initial VA benchmarks across both simulation and real-world environments. To address these tasks, we proposed an explicit Graph-as-Policy (GaP) representation for agentic robot programming, a directed computation graph formulation that encapsulates atomic perception, planning, and control operations into interpretable and modular robot policies. We implemented multi-LLM agent harness that autonomously decomposes natural language specifications to collaboratively generate robot behavior and execution graphs. To standardize and support this generation, we introduced the Modular Open Robot Skill Library (MORSL), an open, evolving framework comprising 51 high-quality, generalizable skill specifications. Furthermore, we presented a simulation-based self-learning mechanism for graph refinement and verification, enabling the system to autonomously debug and tune execution graph structures using internal physics simulations. Finally, we validated our framework with extensive evaluation data, baselines, and ablations on both simulated and real-world physical platforms.

\textbf{Limitations.} Although in experiments, GaP improves success rates significantly over baselines, execution reliability is not yet at industry levels and additional self-learning and parameter tuning is required.  Similarly, GaP execution times are still well below industry standards of 500 units per hour (7 seconds per instance);  more work is required to reduce VLM inference requests and IK motion planning time during execution.  Also, the 8 VA benchmarks focus on quasi-static pick-and-place operations -- only Cable Insertion (IV) requires force sensing.  More work is required to apply GaP to tasks with deformable objects, dynamic forces, and moving targets.  

%provides a robust architecture for composing behaviors, overall system performance remains bounded by the inherent reliability of the underlying skills. Furthermore, our simulation-based graph refinement currently struggles with skills exhibiting a substantial sim-to-real gap, such as end-to-end visuomotor policies. Future work will explore methods to integrate and reliably evaluate these  policies within the graph framework.